\myChapter{État de l'art sur la création d'applications web pilotées par des modèles BPMN}{État de l'art}
\label{chap:etat-art-bpmn-framework-web}
\myMiniToc{mysection}{Sommaire}

\providecommand{\includeChapterImage}[2]{%
  \IfFileExists{#1}{%
    \includegraphics[width=#2]{#1}%
  }{%
    \fbox{\parbox[c][3cm][c]{0.8\textwidth}{%
      \centering Illustration à ajouter : \texttt{\detokenize{#1}}%
    }}%
  }%
}

\mySection{Introduction}{}

Le chapitre précédent a montré qu'un modèle de processus métier peut
servir de pont entre le fonctionnement réel d'une organisation et la
conception d'une application web. Cette relation est particulièrement
importante lorsque le processus est exprimé avec BPMN (Business Process
Model and Notation), car cette notation est compréhensible par les
acteurs métier tout en étant suffisamment structurée pour être exploitée
par des outils logiciels \cite{omg2013bpmn}. Toutefois, représenter un processus ne
signifie pas encore produire une application web complète. Entre le
diagramme BPMN et le code source d'une application, il existe plusieurs
niveaux intermédiaires : les données, les interfaces, les services, les
règles d'accès, les tests, le déploiement et la maintenance.

La littérature montre que plusieurs familles d'approches tentent de
réduire cet écart. Les systèmes de gestion de processus métier, ou BPMS
(Business Process Management Systems), permettent de modéliser, exécuter
et surveiller des processus. L'ingénierie dirigée par les modèles, ou
MDE (Model-Driven Engineering), propose de transformer progressivement
des modèles abstraits vers des artefacts logiciels. Les plateformes
low-code et no-code cherchent, quant à elles, à accélérer la création
d'applications à partir de représentations visuelles. Ces approches sont
proches du sujet de ce mémoire, mais elles ne répondent pas toutes au
même besoin.

Ce chapitre présente donc un état de l'art centré sur une question
précise : comment les travaux existants permettent-ils de passer de
modèles BPMN vers des artefacts d'applications web, et quelles limites
subsistent pour construire un framework capable de générer une
application web maintenable, évolutive et maîtrisable dans son code
source ? Pour éviter une présentation trop générale, l'analyse est
organisée autour d'une démarche de revue, de critères de comparaison
explicites, d'une lecture critique des travaux retenus, puis d'un
positionnement de notre recherche.

\mySection{Démarche de revue et stratégie de recherche documentaire}{}

\subsection{Démarche adoptée}

La revue de littérature a été conduite dans une logique systématique,
sans prétendre reproduire intégralement toutes les étapes d'une revue
systématique formelle. L'objectif n'était pas seulement de rassembler
des références sur BPMN, mais de sélectionner les travaux réellement
utiles pour comprendre le passage du modèle de processus vers
l'application web. Cette démarche rejoint l'esprit des revues
systématiques en génie logiciel, où les travaux sont choisis selon leur
pertinence, leur niveau de détail, leur contribution et leur capacité à
éclairer les limites du domaine \cite{kitchenham2007}.

Concrètement, la recherche documentaire a suivi trois moments. Le
premier a consisté à identifier les notions centrales du sujet : BPMN,
BPMS, applications web orientées processus, transformation de modèles,
génération de code, MDE, MDA (Model-Driven Architecture), IFML
(Interaction Flow Modeling Language), API REST (Representational State
Transfer) et plateformes low-code. Le deuxième moment a porté sur la
sélection des travaux académiques et techniques qui proposent une
approche, une architecture, une méthode ou un outil. Le troisième moment
a consisté à comparer ces travaux à partir des mêmes critères, afin de
dégager les manques que notre travail cherche à traiter.

\subsection{Questions de recherche}

La revue est guidée par une question principale :

\begin{quote}
\textbf{RQ0 : Quelles approches permettent de passer de modèles BPMN
vers des applications web ou des artefacts logiciels, et quelles limites
restent ouvertes pour la génération d'applications web complètes ?}
\end{quote}

Cette question est déclinée en questions secondaires, car le problème
étudié n'est pas un problème unique. Générer une application à partir
d'un modèle BPMN exige de comprendre les capacités du modèle, les outils
existants, les niveaux de transformation, la maîtrise du code et
l'évolution de l'application.

\begin{table}[H]
\caption{Questions de recherche et motivations}
\label{tab:questions-recherche-motivations}
\centering
\footnotesize
\renewcommand{\arraystretch}{1.2}
\begin{tabular}{|p{1.2cm}|p{6.0cm}|p{6.0cm}|}
\hline
\textbf{Code} & \textbf{Question de recherche} & \textbf{Motivation} \\
\hline
QR1 & Quel rôle BPMN joue-t-il dans la conception et l'automatisation
des systèmes d'information orientés processus ? & Comprendre si BPMN est
seulement documentaire ou s'il peut être utilisé comme point d'entrée
dans une chaîne de transformation. \\
\hline
QR2 & Que proposent les BPMS et moteurs de workflows BPMN pour
construire des applications métier ? & Distinguer l'exécution d'un
processus de la génération d'une application web complète. \\
\hline
QR3 & Comment les travaux transforment-ils BPMN vers des interfaces, des
API, des services ou du code ? & Identifier les artefacts effectivement
produits et les informations supplémentaires nécessaires. \\
\hline
QR4 & Comment les approches MDE/MDA traitent-elles la génération, la
traçabilité et l'évolution ? & Évaluer si les modèles peuvent servir de
base à une génération progressive et maintenable. \\
\hline
QR5 & Quelles limites persistent pour une approche source-owned ? &
Vérifier si le code produit reste accessible, versionnable, modifiable
et indépendant d'une plateforme fermée. \\
\hline
\end{tabular}
\end{table}

\subsection{Sources consultées et critères de sélection}

Les sources consultées regroupent trois catégories. La première concerne
les documents de normalisation et les ouvrages de référence, notamment
la spécification BPMN de l'Object Management Group (OMG), les travaux de
Dumas et al. sur le Business Process Management \cite{dumas2018bpm}, ainsi que les
références liées à l'architecture du Web et au développement dirigé par
les modèles. La deuxième catégorie concerne les articles scientifiques
publiés dans des revues, actes de conférences ou plateformes
académiques, portant sur BPMN, IFML, MDA, génération de code,
applications web orientées processus et low-code. La troisième catégorie
regroupe les documentations officielles d'outils, en particulier
Camunda, Flowable, Activiti et jBPM, parce que ces outils sont
fréquemment utilisés comme moteurs d'exécution BPMN dans des
environnements réels \cite{camunda2026modeler,flowable2026bpmn,jbpm2026docs,activiti2026docs}.

Les travaux retenus devaient répondre à au moins l'un des critères
suivants : proposer une méthode de transformation de modèles, présenter
une plateforme ou un moteur BPMN, traiter de la génération d'interfaces
ou d'API, discuter la maintenance du code généré, ou analyser les
limites des approches low-code et MDE. Les documents trop généraux, sans
lien direct avec la génération ou l'exécution de processus, ont été
écartés. Cette sélection permet de maintenir le chapitre centré sur le
sujet du mémoire.

\mySection{Critères d'analyse des travaux existants}{}

Pour éviter une revue composée uniquement de résumés d'articles, les
travaux retenus sont comparés à partir de critères identiques. Ces
critères ont été choisis parce qu'ils correspondent directement au
problème étudié : produire une application web à partir d'un modèle
BPMN, sans perdre la maîtrise du code et sans enfermer l'utilisateur
dans une plateforme difficile à faire évoluer.

\begin{table}[H]
\caption{Critères d'analyse des approches étudiées}
\label{tab:criteres-analyse-approches}
\centering
\footnotesize
\renewcommand{\arraystretch}{1.2}
\begin{tabular}{|p{4.2cm}|p{9.0cm}|}
\hline
\textbf{Critère} & \textbf{Question posée lors de l'analyse} \\
\hline
Modèle d'entrée & Le travail part-il d'un modèle BPMN standard, d'un
modèle enrichi, d'un diagramme complémentaire ou d'un modèle
propriétaire ? \\
\hline
Artefacts produits & L'approche produit-elle des interfaces, des API,
des services backend, une base de données, des tests, une configuration
de moteur ou seulement une documentation ? \\
\hline
Niveau d'automatisation & La transformation est-elle manuelle,
semi-automatique ou automatique ? \\
\hline
Couverture applicative & L'approche couvre-t-elle seulement le workflow
ou bien l'ensemble frontend, backend, données, sécurité et déploiement
? \\
\hline
Maîtrise du code source & Le code est-il lisible, exportable,
versionnable et modifiable hors de la plateforme ? \\
\hline
Évolution et régénération & L'approche prévoit-elle la modification du
modèle, la régénération et la protection des personnalisations ? \\
\hline
\end{tabular}
\end{table}

Ces critères orientent toute la suite du chapitre. Ils permettent de
distinguer clairement une solution qui exécute un processus, une
solution qui assiste la conception, une solution qui génère une
interface, et une solution qui produit réellement une application web
complète.

\mySection{Revue de la littérature}{}

\subsection{BPMN comme point d'entrée pour les systèmes orientés processus}

BPMN est devenu l'un des standards les plus utilisés pour représenter
les processus métier. Selon l'OMG (Object Management Group), BPMN vise à
fournir une notation compréhensible par les parties prenantes qui
conçoivent, pilotent et réalisent les processus, tout en restant assez
précise pour être traduite vers des composants logiciels \cite{omg2013bpmn}. Cette
double vocation explique son intérêt dans un mémoire portant sur les
applications web pilotées par des modèles : BPMN n'est pas seulement un
dessin ; il peut devenir une source de structuration pour l'application.

\begin{figure}[H]
\centering
\includeChapterImage{Chap2/images/image1.png}{0.65\textwidth}
\caption{Extrait des familles d'éléments BPMN : événements, activités,
passerelles, flux et données. Source : OMG, BPMN 2.0.2 \cite{omg2013bpmn}.}
\label{fig:elements-bpmn}
\end{figure}

La figure \ref{fig:elements-bpmn} montre que BPMN organise la description d'un processus
autour d'éléments qui ont une correspondance possible dans une
application web. Une activité utilisateur peut devenir un écran de
traitement ; une tâche de service peut devenir une opération backend ;
une passerelle peut être traduite en règle de décision ; un couloir peut
inspirer une politique de rôles et d'autorisations. Cette correspondance
reste cependant une interprétation. BPMN ne décrit pas directement les
champs d'un formulaire, les tables de base de données, les règles de
validation, les routes d'une API ou les mécanismes d'authentification.

Les travaux sur la qualité des modèles de processus insistent sur cette
difficulté. Un modèle BPMN mal structuré peut être acceptable pour une
discussion métier, mais difficilement exploitable par un générateur. Des
problèmes tels que des activités mal nommées, des passerelles ambiguës,
des événements non reliés ou une mauvaise séparation des rôles peuvent
entraîner des erreurs dans les artefacts générés. Ainsi, avant toute
transformation, le modèle doit être validé sur le plan syntaxique, mais
aussi sur le plan sémantique : chaque activité doit avoir un rôle,
chaque décision doit avoir une condition compréhensible, et chaque
donnée importante doit être identifiable.

Dans cette perspective, BPMN constitue un bon point d'entrée, mais pas
un modèle complet d'application. Il formalise le comportement métier,
c'est-à-dire l'enchaînement des actions et des décisions. Pour générer
une application web, il doit être enrichi par d'autres informations :
données manipulées, acteurs, formulaires, droits d'accès, contraintes de
sécurité et choix techniques. Cette limite n'annule pas l'intérêt de
BPMN ; elle montre plutôt la nécessité d'une chaîne de transformation
plus complète.

\subsection{BPMS et moteurs de workflows BPMN}

Les BPMS, ou Business Process Management Systems, sont les outils les
plus directement associés à BPMN. Ils permettent généralement de
concevoir un processus, de le déployer sur un moteur, d'exécuter des
instances, d'affecter des tâches, de suivre l'état d'avancement et de
collecter des informations de supervision. Ils répondent donc à un
besoin essentiel : faire vivre le processus dans un environnement
informatique.

Camunda constitue l'un des exemples les plus connus dans ce domaine. Son
écosystème propose des outils de modélisation, d'exécution, de gestion
de tâches, d'intégration et de supervision. La documentation de Camunda
indique que le Modeler permet de créer des diagrammes BPMN exécutables
et que le principe recherché est que le diagramme conçu soit celui qui
s'exécute \cite{camunda2026modeler}. Cette idée est importante, car elle limite l'écart
entre la documentation du processus et son comportement réel.

\begin{figure}[H]
\centering
\includeChapterImage{Chap2/images/image2.png}{0.8\textwidth}
\caption{Exemple d'environnement de modélisation BPMN dans Camunda
Modeler. Source : documentation Camunda \cite{camunda2026modeler}.}
\label{fig:camunda-modeler}
\end{figure}

La figure \ref{fig:camunda-modeler} illustre l'intérêt pratique d'un outil comme Camunda : le
concepteur peut représenter visuellement le processus, configurer
certaines tâches et préparer son exécution. Toutefois, même dans un tel
environnement, l'application web finale n'est pas automatiquement
complète. Le moteur sait exécuter le workflow, mais les interfaces
métier, les règles fines de validation, les pages personnalisées, les
modèles de données et l'organisation du code restent souvent à
construire ou à adapter.

Flowable suit une logique proche. Sa documentation le présente comme un
moteur BPMN léger écrit en Java, capable de déployer des définitions de
processus BPMN 2.0, de créer des instances, d'exécuter des requêtes et
d'accéder aux instances actives ou historiques \cite{flowable2026bpmn}. jBPM est
également un moteur de workflow extensible, écrit en Java, orienté vers
l'exécution de processus BPMN 2.0 et intégrable dans des applications
métiers \cite{jbpm2026docs}. Activiti appartient à la même famille d'outils et a
longtemps servi de base à plusieurs solutions de workflow open source.

Ces moteurs sont utiles, mais ils doivent être analysés avec prudence
par rapport au sujet de ce mémoire. Leur objectif principal est
l'orchestration et l'exécution des processus. Or, une application web
métier ne se limite pas à un moteur de workflow. Elle doit aussi gérer
l'expérience utilisateur, les profils, les écrans, les données, les
notifications, les tests, la journalisation, les règles d'accès et le
déploiement. La différence entre un moteur BPMN et un framework de
génération d'applications web est donc fondamentale : le premier exécute
le processus ; le second doit produire une architecture logicielle
exploitable autour du processus.

Cette distinction permet de comprendre pourquoi les BPMS classiques ne
suffisent pas toujours pour une approche dite source-owned. Dans une
approche source-owned, l'équipe de développement doit rester
propriétaire du code source produit : elle doit pouvoir l'ouvrir, le
versionner, le modifier, le tester et le déployer sans dépendre
entièrement d'un environnement propriétaire. Les BPMS peuvent être
intégrés dans une application, mais ils ne garantissent pas à eux seuls
cette maîtrise globale du projet logiciel.

\subsection{Transformation de BPMN vers interfaces, API et artefacts web}

Une deuxième famille de travaux cherche à transformer BPMN vers des
artefacts plus proches du développement web. Ces travaux sont
particulièrement importants pour notre sujet, car ils ne s'arrêtent pas
à l'exécution du processus. Ils posent la question suivante : que
peut-on déduire d'un modèle BPMN pour construire une interface, une API
ou une partie de l'application ?

Les travaux de Sajji, Rhazali et Hadi proposent une transformation
semi-automatique de BPMN vers IFML dans une démarche MDA (Model-Driven
Architecture) \cite{sajji2022}. L'idée est de partir d'un modèle de processus au
niveau CIM (Computation Independent Model), puis de produire un modèle
d'interaction au niveau PIM (Platform Independent Model). IFML, ou
Interaction Flow Modeling Language, sert alors à représenter les vues,
les événements utilisateur, les flux de navigation et les interactions
d'une application web. Cette approche montre clairement que BPMN peut
fournir des informations utiles pour générer des interfaces, mais elle
confirme aussi que BPMN doit être complété par un modèle d'interaction.

Dans un travail complémentaire, les mêmes auteurs étendent la
transformation jusqu'au niveau PSM (Platform Specific Model),
c'est-à-dire vers une représentation plus proche d'une plateforme cible
\cite{sajji2023}. Cette progression est intéressante, car elle montre qu'une
génération réaliste ne se fait pas en une seule étape. Elle passe par
des niveaux successifs : compréhension métier, modèle d'interaction,
modèle technique, puis code. Pour notre recherche, cette idée est
essentielle : un framework de génération à partir de BPMN doit éviter la
traduction directe et fragile du diagramme vers le code. Il doit
introduire des artefacts intermédiaires capables de stabiliser
l'information avant la génération.

D'autres travaux portent sur la génération d'API autour des moteurs de
workflow. Pietruck et Steffens proposent une méthode pour générer
automatiquement des API REST spécifiques au domaine à partir de modèles
de processus et de données \cite{pietruck2020}. Leur contribution répond à un
problème très concret : les moteurs de workflow exposent souvent des API
techniques centrées sur les instances, les tâches et les variables du
moteur, alors que les développeurs d'applications web manipulent plutôt
des objets métier comme demande, dossier, validation, document ou
notification. Une API orientée domaine permet donc de rapprocher le
workflow des concepts compréhensibles par l'application.

Ces travaux montrent que la transformation de BPMN vers le Web peut
prendre plusieurs directions : génération d'interfaces, génération
d'API, configuration d'un moteur, génération de services ou création de
prototypes. Cependant, les approches restent souvent partielles.
Certaines couvrent l'interface, mais pas le backend. D'autres couvrent
les API, mais pas les écrans. D'autres encore supposent l'existence d'un
moteur de workflow. La question d'une génération cohérente, couvrant à
la fois frontend, backend, données, sécurité et tests, demeure donc
ouverte.

Un travail particulièrement proche de cette orientation est celui de
Hoyos et Restrepo-Calle, qui proposent une génération automatique de
prototypes d'applications web orientées processus à partir d'une
spécification en langage naturel restreint \cite{hoyos2017}. Leur idée est de ne
pas demander au concepteur de coder directement l'application, mais de
décrire, dans une grammaire contrôlée, les entités du domaine, les
relations entre ces entités, les tâches du processus et les opérations
attendues sur les données. Cette spécification est ensuite analysée par
un lexer et un parser afin de produire les classes du domaine, les
formulaires, les vues et les contrôleurs d'une application web.

L'intérêt de cette approche, pour notre thème, est qu'elle relie
explicitement trois niveaux souvent traités séparément : les exigences,
les modèles de processus et le code d'une application web. Les auteurs
s'appuient sur des correspondances entre un modèle entité-association,
un modèle BPMN et une représentation textuelle contrôlée. Le prototype
généré peut être exécuté rapidement afin d'être présenté aux parties
prenantes pendant une séance de validation. Cette logique de prototypage
rapide montre que la génération peut être utilisée non seulement pour
produire du code, mais aussi pour réduire les erreurs de compréhension
au début du cycle de développement \cite{hoyos2017}.

\begin{figure}[H]
\centering
\includeChapterImage{Chap2/images/image4.png}{0.65\textwidth}
\caption{Processus de génération de code à partir d'une spécification
textuelle restreinte. Source : Hoyos et Restrepo-Calle \cite{hoyos2017}.}
\label{fig:generation-specification-textuelle}
\end{figure}

Cependant, la proposition reste surtout orientée vers le prototypage.
Les interfaces produites sont fonctionnelles mais limitées, certaines
adaptations visuelles ou règles métier avancées nécessitent encore une
intervention manuelle, et le langage restreint doit être appris par le
concepteur. Cette limite est importante pour notre recherche : elle
montre qu'un framework fondé sur BPMN doit prévoir à la fois une
génération automatique et des mécanismes de personnalisation contrôlée,
afin que l'application obtenue puisse dépasser le simple prototype.

\subsection{MDE, MDA et génération progressive de code}

L'ingénierie dirigée par les modèles, ou MDE (Model-Driven Engineering),
constitue le cadre théorique le plus adapté pour comprendre les
transformations successives nécessaires dans ce mémoire. Dans cette
approche, les modèles ne sont pas de simples documents d'accompagnement.
Ils deviennent des artefacts centraux du développement : ils peuvent
être analysés, validés, transformés et utilisés pour produire du code ou
des configurations.

\begin{figure}[H]
\centering
\includeChapterImage{Chap2/images/image3.png}{0.9\textwidth}
\caption{Principe de transformation dans une approche MDA. Source :
OMG, MDA Guide \cite{omg2014mda}.}
\label{fig:transformation-mda}
\end{figure}

La figure \ref{fig:transformation-mda} illustre l'idée principale de la MDA : passer d'un niveau
abstrait vers un niveau plus concret au moyen de transformations. Dans
le cas de ce mémoire, BPMN peut être vu comme un modèle proche du niveau
métier. Il décrit le déroulement du processus, mais pas encore
l'architecture technique de l'application. Pour générer une application
web, il faut donc produire des modèles ou contrats intermédiaires
décrivant les actions, les rôles, les entités, les écrans, les
transitions, les règles d'accès et les services.

Cette transformation progressive présente plusieurs avantages.
Premièrement, elle rend la génération plus compréhensible : chaque étape
produit un artefact qui peut être inspecté. Deuxièmement, elle facilite
les tests : il devient possible de vérifier le modèle BPMN, puis les
contrats intermédiaires, puis les fichiers générés. Troisièmement, elle
limite le couplage entre BPMN et un framework web particulier. Si l'on
change de technologie cible, par exemple de Laravel à Spring Boot ou de
Vue.js à Angular, les contrats intermédiaires peuvent rester stables
tandis que seuls les générateurs de code changent.

Les travaux sur l'intégration du code généré et du code écrit
manuellement montrent également une difficulté importante
\cite{greifenberg2015,rahad2021}. Dans les
projets réels, le développeur doit souvent personnaliser ce qui est
produit automatiquement : ajouter une règle métier, modifier une
interface, brancher un service externe, ajuster une requête ou corriger
un comportement spécifique. Si le générateur écrase ces modifications
lors d'une nouvelle génération, il devient dangereux. S'il ne permet
jamais de régénérer, il perd son intérêt. Une approche robuste doit donc
prévoir la séparation entre fichiers générés, fichiers personnalisables
et extensions manuelles.

Cette question est centrale pour un framework de création d'applications
web pilotées par BPMN. La génération initiale n'est pas suffisante. Les
processus métier évoluent, les formulaires changent, les rôles se
précisent, et les règles d'organisation sont modifiées. Le framework
doit donc intégrer la régénération contrôlée : détecter les changements
du modèle, identifier les fichiers concernés, préserver les
personnalisations et signaler les conflits éventuels.

Les travaux de Drave et al. s'inscrivent précisément dans cette
problématique. Les auteurs proposent une approche d'ingénierie dirigée
par les modèles pour faire évoluer un système d'information centré sur
les données vers un système d'information conscient des processus, ou
PAIS (Process-Aware Information System) \cite{drave2022}. Leur démarche combine
plusieurs langages spécifiques au domaine : modèles de données,
contraintes OCL (Object Constraint Language), modèles d'interfaces,
modèles de marquage des rôles et modèles BPMN. Cette composition de
langages permet de produire une application dans laquelle les processus
ne sont pas simplement documentés, mais intégrés à l'architecture
logicielle générée.

Un apport essentiel de ce travail est la proposition d'une notation
textuelle de BPMN destinée à la génération de code. Les auteurs montrent
que la notation graphique BPMN est lisible pour les acteurs métier, mais
qu'une représentation textuelle facilite l'intégration avec les outils
de développement, la gestion de versions et la régénération continue.
Ils ajoutent également une étape de validation des modèles, fondée sur
des conditions de contexte et sur une transformation vers des réseaux de
Petri pour vérifier la correction comportementale du processus \cite{drave2022}.
Cette dimension est importante : dans une approche générative, un modèle
BPMN incorrect peut produire automatiquement une application incorrecte.

\begin{figure}[H]
\centering
\includeChapterImage{Chap2/images/image5.png}{0.9\textwidth}
\caption{Génération en deux étapes d'un système d'information conscient
des processus. Source : Drave et al. \cite{drave2022}.}
\label{fig:generation-pais-deux-etapes}
\end{figure}

La figure \ref{fig:generation-pais-deux-etapes} illustre la logique proposée : le modèle de domaine, le
modèle BPMN et le modèle de marquage sont d'abord transformés en modèles
intermédiaires de données et d'interfaces. Ces modèles sont ensuite
exploités par un générateur pour produire les composants frontend,
backend et les descriptions BPMN exécutables. Cette organisation est
intéressante pour notre mémoire, car elle montre qu'il est préférable de
ne pas passer directement de BPMN vers le code final. Un passage par des
modèles intermédiaires rend la transformation plus traçable et facilite
la maintenance.

La même étude précise toutefois que tout ne peut pas être généré. La
logique métier spécifique, certaines connexions avec les données
persistantes et certaines adaptations de l'interface doivent rester
ouvertes au code écrit manuellement. Les auteurs défendent donc une
génération continue, mais extensible, où les artefacts générés peuvent
cohabiter avec des ajouts humains. Cette idée rejoint directement notre
positionnement : un framework de création d'applications web pilotées
par BPMN doit produire une base cohérente, tout en préservant
l'intervention du développeur lorsque le métier devient trop spécifique.

\subsection{Low-code, no-code et applications orientées processus}

Les plateformes low-code et no-code constituent une autre famille proche
du sujet. Elles proposent de créer rapidement des applications à partir
de modèles visuels, de composants prédéfinis et de connecteurs. Elles
ont l'avantage de réduire le temps de développement et d'impliquer plus
facilement les utilisateurs métier. Plusieurs travaux récents montrent
que ces plateformes partagent certains objectifs avec l'ingénierie
dirigée par les modèles : abstraction, automatisation, réduction du
codage répétitif et accélération du prototypage \cite{diruscio2022,bucaioni2022}.

Cependant, leur logique n'est pas toujours compatible avec une recherche
orientée génération de code maîtrisable. Certaines plateformes enferment
l'application dans un environnement propriétaire. D'autres permettent
l'export du code, mais ce code n'est pas toujours lisible ou facile à
maintenir. Il existe aussi des difficultés de gouvernance :
multiplication d'applications créées rapidement, dépendances implicites,
règles métier dispersées, faible traçabilité et contrôle limité sur
l'architecture.

Les plateformes orientées processus, comme AKIP ou certaines solutions
proches des BPMS, montrent qu'il existe un besoin réel pour des
applications web centrées sur les tâches, les rôles et les instances de
processus \cite{neto2022}. Elles facilitent la création d'applications où le
workflow est au cœur du fonctionnement. Toutefois, elles restent souvent
liées à leur propre environnement d'exécution. Pour notre mémoire,
l'enjeu est différent : il ne s'agit pas seulement de fournir une
plateforme d'exécution, mais de proposer une architecture et des
méthodes permettant de générer un projet web structuré, extensible et
exploitable par une équipe de développement.

La plateforme AKIP proposée par Neto et al. est l'un des travaux les
plus proches de la problématique de ce mémoire, car elle vise
explicitement le développement d'applications web conscientes des
processus à partir de modèles BPMN exécutables \cite{neto2022}. Les auteurs
partent du constat que les plateformes classiques d'automatisation de
processus souffrent de plusieurs limites : dépendance à un fournisseur,
interfaces utilisateur peu flexibles, intégration coûteuse avec d'autres
systèmes, technologies parfois vieillissantes et faible prise en charge
des fonctionnalités non directement liées au processus, comme la gestion
des utilisateurs, les opérations CRUD (Create, Read, Update, Delete) sur
les entités métier ou l'import/export de données.

AKIP répond à ce problème en proposant une architecture de référence et
un générateur d'applications, appelées KIPApps. Une KIPApp est une
application web moderne qui s'exécute au-dessus d'un moteur de processus
et qui fournit, en plus de l'exécution du modèle BPMN, des
fonctionnalités de gestion d'application, de gestion de processus, de
liste de tâches et d'intégration avec des systèmes externes. L'approche
utilise notamment des métadonnées JSON (JavaScript Object Notation) pour
décrire les entités de domaine, les formulaires de démarrage de
processus et les formulaires associés aux tâches utilisateur \cite{neto2022}.

\begin{figure}[H]
\centering
\includeChapterImage{Chap2/images/image6.png}{0.85\textwidth}
\caption{Vue générale de la plateforme AKIP pour la génération
d'applications web conscientes des processus. Source : Neto et al.
\cite{neto2022}.}
\label{fig:plateforme-akip}
\end{figure}

La figure \ref{fig:plateforme-akip} montre que la plateforme repose sur deux éléments
complémentaires : un générateur d'application et une architecture de
référence. Cette organisation est importante, car elle évite de réduire
la génération à la production de quelques fichiers isolés. Elle inscrit
la génération dans un cadre applicatif complet comprenant frontend,
backend, moteur de processus, connecteurs et fonctionnalités natives.

La validation présentée par Neto et al. montre que la plateforme a été
utilisée dans un contexte industriel pour automatiser trois processus
issus du domaine logistique et portuaire \cite{neto2022}. Les résultats
indiquent que les applications générées sont utilisables, mais que des
personnalisations restent nécessaires pour intégrer les règles métier et
adapter certaines interfaces. Les limites signalées par les auteurs
concernent aussi la couverture partielle de certains éléments BPMN,
notamment les types de passerelles et d'événements. Pour notre
recherche, AKIP montre donc une voie concrète : générer une application
web complète à partir de BPMN est possible, mais cette génération doit
être accompagnée d'un modèle de domaine, de métadonnées de formulaires,
d'un mécanisme d'extension et d'une stratégie claire de
personnalisation.

Le low-code apporte donc une inspiration, mais pas une solution
suffisante. Il montre que la modélisation visuelle peut accélérer la
création d'applications. Mais il laisse ouverte la question de la
propriété du code, de la régénération contrôlée et de l'intégration avec
les pratiques classiques du génie logiciel, comme le versionnement, les
tests automatisés et le déploiement continu.

\mySection{Discussion et synthèse critique}{}

L'analyse des travaux montre que le domaine est riche, mais fragmenté.
BPMN apporte une notation standard et compréhensible. Les BPMS apportent
l'exécution et la supervision. Les approches BPMN vers IFML rapprochent
le processus des interfaces. Les travaux sur les API de workflows
rapprochent les moteurs BPMN des applications clientes. La MDE fournit
le cadre méthodologique des transformations. Enfin, le low-code montre
l'intérêt pratique des approches visuelles. Pourtant, aucune de ces
familles ne couvre entièrement le besoin visé dans ce mémoire.

\begin{table}[H]
\caption{Synthèse critique des familles de travaux}
\label{tab:synthese-critique-familles}
\centering
\footnotesize
\renewcommand{\arraystretch}{1.2}
\begin{tabular}{|p{3.3cm}|p{4.8cm}|p{5.2cm}|}
\hline
\textbf{Famille de travaux} & \textbf{Apport principal} & \textbf{Limite pour notre sujet} \\
\hline
BPMN comme modèle métier & Langage standard, lisible par les acteurs
métier, utile pour représenter activités, événements, rôles et
décisions. & Ne décrit pas suffisamment données, formulaires, sécurité,
API et architecture web. \\
\hline
BPMS et moteurs BPMN & Exécution, suivi des instances, gestion des
tâches, intégration avec des systèmes externes. & Centrés sur le
workflow ; ne génèrent pas toujours une application web complète et
source-owned. \\
\hline
BPMN vers IFML / interfaces & Possibilité de dériver des modèles
d'interaction et des interfaces à partir du processus. & Couverture
souvent limitée au frontend ou à la navigation. \\
\hline
API et services de workflow & Production d'API orientées domaine autour
des moteurs de processus. & Ne couvre pas à elle seule les écrans,
données, sécurité, tests et déploiement. \\
\hline
MDE/MDA et génération de code & Cadre solide pour organiser les
transformations successives et la traçabilité. & Nécessite des
mécanismes précis pour gérer le code manuel et la régénération. \\
\hline
Low-code/no-code & Rapidité de prototypage, composants visuels,
implication des utilisateurs métier. & Risque de dépendance à la
plateforme et faible maîtrise du code source. \\
\hline
\end{tabular}
\end{table}

La principale conclusion est que la difficulté ne réside pas seulement
dans la génération de quelques fichiers. Il est relativement possible de
produire une interface simple ou une configuration de moteur à partir
d'un modèle. La difficulté scientifique et technique est de produire une
application web cohérente, structurée en couches, capable d'évoluer et
dont le code reste compréhensible. C'est précisément sur ce point que
les travaux existants restent insuffisants.

Les trois articles ajoutés renforcent cette conclusion. Hoyos et
Restrepo-Calle montrent qu'il est possible de générer rapidement un
prototype web à partir d'une spécification contrôlée combinant entités
et processus \cite{hoyos2017}. Drave et al. montrent qu'une application
consciente des processus peut être générée de manière plus robuste
lorsque BPMN est intégré dans une chaîne MDE avec modèles de données,
interfaces, rôles et validation formelle \cite{drave2022}. Neto et al., à
travers AKIP, montrent enfin qu'une plateforme peut produire des
applications web opérationnelles au-dessus d'un moteur BPMN, mais que
les personnalisations métier et la couverture complète de BPMN restent
des points ouverts \cite{neto2022}. Ainsi, la littérature converge vers une
même idée : BPMN est nécessaire pour piloter le comportement, mais il
doit être complété par des artefacts décrivant les données, les écrans,
les rôles, les services et les règles métier.

Cette convergence éclaire directement le besoin traité dans ce mémoire.
Le problème n'est pas seulement de transformer un diagramme en code,
mais de définir une architecture et une méthode capables d'organiser
tout le passage du modèle de processus vers l'application web. Une telle
méthode doit préciser quels éléments sont extraits de BPMN, quels
éléments doivent être ajoutés dans des modèles complémentaires, comment
la cohérence est vérifiée et comment le code généré peut évoluer sans
perdre les ajouts humains.

Une autre limite concerne la prise en compte de l'évolution. Les
processus métier changent. Une application générée doit donc être
capable de suivre ces changements sans perdre les ajouts manuels
effectués par les développeurs. Or, beaucoup d'approches décrivent la
génération initiale, mais traitent peu la régénération, la comparaison
entre versions de modèles, la gestion des conflits et la préservation
des personnalisations.

Enfin, la littérature montre que BPMN seul ne suffit pas. Il doit être
enrichi ou interprété à travers des contrats intermédiaires. Ces
contrats peuvent décrire les entités, les actions, les rôles, les états,
les écrans, les règles de validation et les services. Ils constituent un
pont entre le modèle métier et l'architecture technique. Cette idée sera
importante dans la suite du mémoire, car elle permet d'éviter une
génération directe, fragile et difficile à maintenir.

\mySection{Positionnement de la recherche}{}

À partir des limites identifiées, notre recherche se positionne à
l'intersection de trois domaines : BPMN comme langage de modélisation
des processus métier, l'ingénierie dirigée par les modèles comme cadre
de transformation, et le développement web comme cible de génération.
Elle ne vise pas à proposer un simple éditeur BPMN, ni à remplacer un
moteur de workflow existant, ni à créer une plateforme low-code fermée.
Elle vise plutôt à définir une architecture et des méthodes pour générer
des applications web pilotées par des modèles BPMN, tout en conservant
la maîtrise du code source.

Le positionnement proposé repose sur cinq principes.

\begin{itemize}
\item
  Utiliser BPMN comme point d'entrée métier, parce qu'il permet de
  représenter les activités, les décisions, les événements et les
  responsabilités dans un langage visuel standard.
\item
  Introduire des contrats intermédiaires afin de compléter BPMN par les
  informations nécessaires à la génération : données, formulaires,
  rôles, transitions, services, validations et règles d'accès.
\item
  Générer des artefacts frontend et backend de manière cohérente, au
  lieu de se limiter à l'interface ou au moteur de workflow.
\item
  Préserver une approche source-owned, où le code produit reste lisible,
  versionnable, modifiable et déployable en dehors d'une plateforme
  fermée.
\item
  Prévoir l'évolution du modèle et la régénération contrôlée, afin de
  protéger les personnalisations humaines et d'accompagner les
  changements métier.
\end{itemize}

Ce positionnement distingue notre contribution des solutions existantes.
Contrairement à un BPMS classique, l'objectif n'est pas seulement
d'exécuter un processus. Contrairement à une plateforme low-code fermée,
l'objectif n'est pas de cacher le code généré. Contrairement aux
approches de prototypage rapide, l'objectif ne se limite pas à produire
une première maquette. Et contrairement aux approches MDE très outillées
comme celles de Drave et al. ou AKIP, notre travail cherche à définir
une architecture et une méthode adaptées à la création progressive
d'applications web pilotées par BPMN, en explicitant les artefacts
intermédiaires nécessaires et les règles de passage vers les composants
applicatifs.

\mySection{Conclusion}{}

Ce chapitre a présenté un état de l'art centré sur la création
d'applications web pilotées par des modèles BPMN. L'analyse a montré que
BPMN constitue un point d'entrée pertinent pour représenter les
processus métier, mais qu'il ne contient pas toutes les informations
nécessaires à la génération d'une application complète. Les BPMS et
moteurs BPMN apportent des capacités solides d'exécution et de
supervision, mais ils restent principalement orientés workflow. Les
approches de transformation vers IFML, API ou code montrent des pistes
intéressantes, mais leur couverture demeure souvent partielle.

La discussion a également mis en évidence l'importance des approches
MDE/MDA pour organiser les transformations et assurer la traçabilité.
Toutefois, la génération de code doit être pensée avec prudence : elle
doit permettre l'intervention humaine, préserver les personnalisations
et accompagner l'évolution du modèle. Les plateformes low-code
confirment l'intérêt des modèles visuels pour accélérer le
développement, mais elles soulèvent des problèmes de dépendance et de
maîtrise du code.

En conséquence, il apparaît nécessaire de proposer une architecture de
framework capable d'exploiter BPMN comme point d'entrée, de produire des
contrats intermédiaires et de générer des artefacts web cohérents. Les
travaux récents analysés dans ce chapitre montrent que cette orientation
est pertinente, mais qu'elle reste ouverte dès que l'on cherche à
concilier génération automatique, qualité des modèles, richesse
fonctionnelle d'une application web, personnalisation métier et
régénération contrôlée. C'est dans cet espace que se situe la
contribution proposée dans la suite du mémoire.

